\documentclass[12pt,a4paper]{article}

\usepackage[utf8]{inputenc}
\usepackage[T1]{fontenc}
\usepackage[portuguese]{babel}

\usepackage{listings}
\usepackage{xcolor}
\usepackage{amsmath}
\usepackage{amssymb}
\usepackage{amsfonts}
\usepackage{graphicx}

\lstset{
    basicstyle=\ttfamily\small,
    keywordstyle=\color{blue}\bfseries, 
    commentstyle=\color{green!70!black},
    stringstyle=\color{red}, 
    numbers=left, 
    numberstyle=\tiny\color{gray}, 
    frame=single, 
    breaklines=true, 
    showstringspaces=false 
}

\usepackage[margin=2.5cm]{geometry}

\begin{document}

\begin{center}
    {\LARGE MC521 - Relatório II} \\ [2,5em]
    {\large Júlio da Silva Telles} \\ [1,0em]
    {\large 17 de Junho de 2026} \\ [1,5em]
\end{center}

\vspace{1cm} 

\section{A - Solve It (UVA - 10341)}

O problema consiste em encontrar a raiz da função contínua $f(x)$, onde:
\begin{center}
$f(x) = p * e^{-x} + q * \sin(x) + r * \cos(x) + s * \tan(x) + t * x^2 + u$
\end{center}

\subsection{Solução}

Como o domínio de $x$ é restrito ao intervalo $[0, 1]$, podemos utilizar o \textbf{Método da Bisseção} (uma adaptação da busca binária para funções contínuas). O algoritmo segue os seguintes passos:

\begin{enumerate}
    \item \textbf{Verificação de Existência:} Pelo Teorema de Bolzano, como a função é contínua, só existirá uma raiz no intervalo se $f(0)$ e $f(1)$ tiverem sinais opostos. Calculamos ambos os valores; se possuírem o mesmo sinal (ambos positivos ou ambos negativos), podemos afirmar com segurança que a equação não possui solução e imprimimos \texttt{"No solution"}.
    
    \item \textbf{Busca da Raiz:} Caso exista uma solução, inicializamos nossos limites como $a = 0$ e $b = 1$. Em cada iteração, calculamos o ponto médio $m = \frac{a + b}{2}$ e avaliamos $f(m)$. Dependendo do sinal de $f(m)$, substituímos o limite $a$ ou $b$ por $m$, reduzindo o intervalo de busca pela metade a cada passo.
    
    \item \textbf{Condição de Parada:} Como estamos lidando com números de ponto flutuante, o algoritmo repete o processo de divisão até que o tamanho do intervalo seja menor que uma tolerância aceitável (neste caso, utilizamos um erro de $\approx 10^{-10}$ para garantir a precisão exigida de 4 casas decimais).
\end{enumerate}

\end{document}